\documentclass[11pt]{article}

\usepackage[margin=1in]{geometry}
\usepackage{amsmath}
\usepackage{amssymb}
\usepackage{booktabs}
\usepackage{xspace}

\newcommand{\modelname}{\textsc{GraphMemory3}\xspace}
\newcommand{\localmem}{\mathcal{M}^{\mathrm{L}}}
\newcommand{\globalmem}{\mathcal{M}^{\mathrm{G}}}
\newcommand{\states}{\mathcal{S}}
\newcommand{\actions}{\mathcal{A}}
\newcommand{\episodes}{\mathcal{D}}
\newcommand{\sig}{\sigma}

\title{\modelname: Abstract and Methodology}
\author{}
\date{}

\begin{document}
\maketitle

\begin{abstract}
Large language model (LLM) agents increasingly operate in interactive environments
where they make sequential decisions under feedback from the environment. While
procedural memory has shown promise for improving long-horizon decision making,
existing memory systems are often centered on individual agents, requiring each
agent to independently collect costly procedural experience through repeated
interaction. This produces sparse, domain-bound, and isolated memories that offer
limited support when agents enter under-explored environments. Recent shared-memory
approaches suggest that experience can transfer across tasks, domains, and agents,
but they mainly emphasize memory representation or storage, leaving open a central
question: how should agents selectively collaborate around memory, deciding what
should remain local, what should be promoted globally, and when each type of memory
should influence action selection?

We propose \modelname, a collaborative episodic graph memory framework for LLM
agents in interactive environments. \modelname converts interaction trajectories
into structured episodic graphs that encode states, actions, outcomes, failures,
and temporal or causal relations. It separates memory into domain-specific local
graph memory and shared global graph memory, promotes transferable procedures into
the global memory, and introduces phase-aware routing to dynamically select memory
roles during interaction. Instead of injecting raw trajectories or all retrieved
memories into the prompt, \modelname composes a concise decision summary from local
grounding, global workflow knowledge, source-role priors, and failure-avoidance
signals. By turning shared experience from passive context into collaborative
decision support, \modelname enables agents to reuse procedural knowledge while
preserving environment-specific grounding. Experiments on interactive
decision-making benchmarks show that \modelname improves task success and reduces
redundant exploration compared with single-agent and shared-memory baselines.
\end{abstract}

\section{Methodology}

\subsection{Overview: Memory-Centric Collaboration}

\modelname is designed around the idea that collaboration between LLM agents
should happen through memory operations rather than through direct action control.
The agent still uses its original solver, prompt template, action parser, and
environment loop. \modelname intervenes only in the memory layer: it decides which
experience should be kept local, which experience should become shared, and how
local and global memory should be routed into the current decision.

This design is motivated by three limitations of existing memory systems. First,
single-agent procedural memories are costly because each agent must rediscover the
same search, manipulation, and recovery patterns. Second, naive shared memory can
be harmful because concrete experience from one environment may be wrong in
another. Third, retrieval alone is insufficient: a useful memory must be selected
according to the current phase and must map to an action that is currently
possible. Thus \modelname implements memory-centric collaboration through four
linked operations:
\begin{equation}
    \tau
    \xrightarrow{\text{construct}}
    G_{\tau}
    \xrightarrow{\text{local update}}
    \localmem_i
    \xrightarrow{\text{promote}}
    \globalmem
    \xrightarrow{\text{route}}
    \mathcal{C}_t .
\end{equation}
Here $\tau$ is an interaction episode, $G_{\tau}$ is its episodic graph,
$\localmem_i$ is local memory for a scene, domain, or agent, $\globalmem$ is shared
global memory, and $\mathcal{C}_t$ is the memory summary injected at decision step
$t$.

\subsection{Interactive Decision Setting}

At step $t$, an agent receives observation $o_t$, goal $g$, and admissible actions
$\actions_t$, then selects $a_t \in \actions_t$. The environment returns feedback
$y_t$, reward or success signal $r_t$, and the next observation $o_{t+1}$:
\begin{equation}
    \tau =
    \bigl\{(o_t, a_t, y_t, r_t, o_{t+1})\bigr\}_{t=1}^{T}.
\end{equation}
The LLM policy is conditioned on a memory context $\mathcal{C}_t$:
\begin{equation}
    a_t \sim
    \pi_{\theta}
    \left(a \mid o_t, g, \actions_t, \mathcal{C}_t\right).
\end{equation}
The key methodological question is how to construct $\mathcal{C}_t$. Instead of
retrieving a top-$k$ list of memories, \modelname treats $\mathcal{C}_t$ as a
routed collaboration product: local memory contributes environment grounding,
global memory contributes transferable procedure, and the router decides which
role should be active in the current phase.

\subsection{Step 1: Construct Episodic Graphs}

\textbf{Why.} Raw trajectories are difficult to share because they mix together
state descriptions, actions, task-specific wording, failures, and incidental
details. Directly sharing them can create prompt noise and poor transfer. We
therefore convert trajectories into episodic graphs, which preserve procedural
structure while allowing repeated patterns to be counted and compared.

\textbf{How.} Each episode is converted into
\begin{equation}
    G_{\tau} = (V_{\tau}, E_{\tau}),
\end{equation}
with node types
\begin{equation}
    \mathcal{T}_V =
    \{\textsc{State}, \textsc{Action}, \textsc{Subgoal}, \textsc{Failure}\},
\end{equation}
and edge types
\begin{equation}
    \mathcal{T}_E =
    \{\textsc{Temporal}, \textsc{Causes}, \textsc{AdvancesTo},
    \textsc{FailsUnder}\}.
\end{equation}
For a transition, \modelname adds the execution chain
\begin{equation}
    v_t^{s} \rightarrow v_t^{a} \rightarrow v_{t+1}^{s},
\end{equation}
and attaches outcome edges to subgoal or failure nodes when the action advances
or blocks progress. Actions are canonicalized into slot-level forms such as
\begin{equation}
    \texttt{take(object=apple, source=countertop)}.
\end{equation}

\textbf{Collaboration role.} This step makes experience shareable. Different
agents may use different language or encounter different object instances, but the
graph abstraction exposes comparable procedural events such as ``search a source,''
``take target,'' ``process held object,'' or ``deliver to destination.''

\subsection{Step 2: Maintain Local Memory for Grounding}

\textbf{Why.} Not all experience should be shared globally. Some memories are
layout-specific, scene-specific, or dependent on the current action frontier. For
example, a cup being found on \texttt{countertop 3} is highly useful in the same
kitchen, but dangerous if transferred as a global fact. \modelname therefore keeps
environment-grounded evidence in local memory.

\textbf{How.} For local context $i$, \modelname maintains
\begin{equation}
    \localmem_i =
    \left(V_i, E_i, \mathcal{P}_i, \mathcal{R}_i, \mathcal{H}_i\right),
\end{equation}
where $V_i,E_i$ are aggregated graph nodes and edges, $\mathcal{P}_i$ are
promotion candidates, $\mathcal{R}_i$ are induced rules, and $\mathcal{H}_i$ are
memory artifacts. Graph elements are merged by signature $\sig(\cdot)$, and each
element $x$ stores support, positive evidence, negative evidence, and stalled
evidence:
\begin{equation}
    \mathrm{stat}(x) = (n_x,p_x,q_x,z_x).
\end{equation}
The confidence and stall rate are
\begin{equation}
    \mathrm{conf}(x)=\frac{p_x}{p_x+q_x+\epsilon},
    \qquad
    \mathrm{stall}(x)=\frac{z_x}{n_x+\epsilon}.
\end{equation}

\textbf{Collaboration role.} Local memory protects collaboration from
over-generalization. It lets an agent exploit precise scene evidence while keeping
that evidence separate from global procedural knowledge.

\subsection{Step 3: Induce Procedural Units}

\textbf{Why.} A graph memory is useful only if it can be converted into decision
support. Storing every state and edge is too low-level for LLM prompting, while
sharing only natural-language reflections loses structure. \modelname therefore
induces intermediate procedural units that can be retrieved, promoted, and routed.

\textbf{How.} From the local graph, \modelname induces candidates for
preconditions, failure patterns, workflow chunks, and repairs:
\begin{equation}
    \mathcal{P}_i =
    \mathcal{P}^{\mathrm{pre}}_i
    \cup \mathcal{P}^{\mathrm{fail}}_i
    \cup \mathcal{P}^{\mathrm{work}}_i
    \cup \mathcal{P}^{\mathrm{repair}}_i .
\end{equation}
These candidates are compressed into rules and artifacts. Rules describe
phase-conditioned procedural knowledge, while artifacts expose reusable prompt
units such as plans, scene relations, source-type priors, and reflections. Two
artifact types are central:
\begin{align}
    h_{\mathrm{rel}} &: \text{target object} \rightarrow \text{local source instance},
    \\
    h_{\mathrm{prior}} &: \text{task or target family} \rightarrow \text{source type}.
\end{align}

\textbf{Collaboration role.} This step creates the objects over which agents
collaborate. Rather than sharing whole episodes, agents share compact procedural
claims with statistics, provenance, and scope.

\subsection{Step 4: Promote Transferable Memory Globally}

\textbf{Why.} Shared memory should contain procedures that are likely to transfer,
not arbitrary details. A global memory that stores concrete locations can mislead
agents in new environments. Conversely, keeping all memories local prevents agents
from benefiting from each other's costly exploration. \modelname promotes only
items whose evidence suggests reusable procedural value.

\textbf{How.} The global memory is
\begin{equation}
    \globalmem =
    \left(
    \mathcal{P}^{\mathrm{G}},
    \mathcal{R}^{\mathrm{G}},
    \mathcal{H}^{\mathrm{G}}
    \right).
\end{equation}
For a rule or artifact $x$, promotion depends on support, confidence, coverage,
utility, transfer success, specificity, conflict, and stalled evidence. A rule
utility can be written as
\begin{equation}
\begin{aligned}
    U(x) ={}&
    1.1\,\mathrm{conf}(x)
    + 0.45\,\min\!\left(\frac{\mathrm{supp}(x)}{4},1\right)
    + 0.35\,\min\!\left(\frac{\mathrm{cov}(x)}{2},1\right) \\
    &+ 0.40\,\mathrm{util}(x)
    + 0.25\,\mathrm{trans}(x)
    - 0.35\,\mathrm{spec}(x)
    - 0.30\,\mathrm{conflict}(x)
    - 0.45\,\mathrm{stall}(x).
\end{aligned}
\end{equation}
An item is promoted when
\begin{equation}
    x \in \globalmem
    \quad \Longleftarrow \quad
    U(x) \ge \lambda_{\mathrm{promote}}
    \;\wedge\;
    \mathrm{cov}(x) \ge \lambda_{\mathrm{cov}}.
\end{equation}
In practice, concrete scene relations remain local, while source-type priors,
workflow chunks, and repair patterns are more suitable for global memory.

\textbf{Collaboration role.} Promotion is the collaboration boundary. Local agents
contribute experience to a shared pool only when the evidence indicates that the
memory is reusable beyond its original scene.

\subsection{Step 5: Retrieve Local and Global Evidence Separately}

\textbf{Why.} Local and global memory serve different purposes. Local memory should
answer ``what is grounded in this environment now?'' Global memory should answer
``what procedure tends to work across environments?'' Mixing them in a single
retrieval list hides this distinction.

\textbf{How.} At step $t$, \modelname builds a structured query
\begin{equation}
    q_t =
    \left(g, s_t, \phi_t, \rho_g, \actions_t, b_t\right),
\end{equation}
where $s_t$ is scene or domain context, $\phi_t$ is the progress phase, $\rho_g$
contains goal roles such as target, tool, and destination, and $b_t$ records
dynamic belief such as visible objects, held objects, searched locations, and
exhausted sources. Candidate memory items are scored by
\begin{equation}
    \mathrm{score}(m,q_t)
    =
    \alpha_s R_s(m,q_t)
    + \alpha_{\tau} R_{\tau}(m,q_t)
    + \alpha_g R_g(m,q_t)
    + \alpha_c \mathrm{conf}(m).
\end{equation}
The retrieval output keeps local and global evidence in separate roles, rather
than collapsing them into one undifferentiated memory context.

\textbf{Collaboration role.} This separation lets the current agent combine its
own grounded experience with collaborators' transferable experience without
confusing the authority of the two memory sources.

\subsection{Step 6: Arbitrate Source Evidence}

\textbf{Why.} In search tasks, many failures come from repeatedly exploring poor
sources. However, source priors can also be noisy: one agent's successful source
type may not be available, may already be exhausted, or may only weakly match the
current target. \modelname therefore does not directly prompt every source prior.
It first builds a source evidence table that arbitrates positive, negative, local,
and global evidence.

\textbf{How.} For a source base $b$, evidence is aggregated from local and global
source-type priors, scene relations, searched-empty relations, current exhausted
locations, and current admissible actions. With default weights
\begin{equation}
    w_{\mathrm{local}}=1.0,
    \qquad
    w_{\mathrm{global}}=0.65,
\end{equation}
positive evidence $e$ contributes
\begin{equation}
    S^+_e =
    w_{m(e)}
    \left(
    0.7\,\mathrm{conf}(e)
    + 0.3\,\ell(n_e)
    + 0.04\,\max(0,p_e-q_e)
    + \beta_e
    \right)
    \gamma_e ,
\end{equation}
where
\begin{equation}
    \ell(n_e)=\frac{1}{1+\exp[-0.35(n_e-2)]},
    \qquad
    \gamma_e =
    \begin{cases}
    1.00, & \text{exact target match},\\
    0.48, & \text{same task family},\\
    0.25, & \text{role-only evidence}.
    \end{cases}
\end{equation}
The final source score is
\begin{equation}
    F(b)=
    \sum_{e\rightarrow b} S^+_e
    -
    \sum_{e\rightarrow b} S^-_e
    -
    0.20\,\min(E_b,6),
\end{equation}
where $E_b$ is the number of times source base $b$ has been exhausted. A source
can become a next priority only if it has exact or same-family evidence, passes a
score threshold, and maps to an admissible action.

\textbf{Collaboration role.} This table is the operational mechanism that turns
shared search experience into collaborative decision support. It lets global
memory bias exploration, while local memory and current feedback can veto or
ground that bias.

\subsection{Step 7: Route Memory by Decision Phase}

\textbf{Why.} The same memory can be useful in one phase and harmful in another.
For example, a source prior is useful while searching for a target, but it is
noise after the target is held. A global workflow can guide high-level progress,
but should not override a concrete local action. \modelname therefore routes
memory according to phase and role.

\textbf{How.} \modelname constructs five section types:
\begin{equation}
    Z_t =
    \{
    z_{\mathrm{phase}},
    z_{\mathrm{local}},
    z_{\mathrm{source}},
    z_{\mathrm{global}},
    z_{\mathrm{failure}}
    \}.
\end{equation}
Each section is assigned a lightweight text loss:
\begin{equation}
    \mathcal{L}(z,q_t)
    =
    1
    -G_{\mathrm{phase}}
    -G_{\mathrm{transfer}}
    -G_{\mathrm{ground}}
    -G_{\mathrm{act}}
    +P_{\mathrm{noise}}
    +P_{\mathrm{wrong}}.
\end{equation}
The router then composes memory by role:
\begin{equation}
    Z_t^\star =
    \mathrm{Compose}
    \left(
    z_{\mathrm{phase}},
    z_{\mathrm{local}},
    z_{\mathrm{source}},
    z_{\mathrm{global}},
    z_{\mathrm{failure}};
    \phi_t
    \right).
\end{equation}
Phase policy anchors the prompt, local grounding bridges memory to the current
state, source roles are emphasized in search phases, global workflow remains
advisory, and failure memory suppresses repeated mistakes.

\textbf{Collaboration role.} Routing decides when each memory contributor should
speak. This is what makes the system collaborative rather than merely shared: the
agent is not given a bag of memories, but a phase-specific composition of local
and global roles.

\subsection{Step 8: Emit Only Actionable Memory}

\textbf{Why.} Extra memory can degrade LLM decisions when it is weak, stale, or not
actionable. In particular, a global prior such as ``countertops are often useful''
should not become a command unless a corresponding admissible action exists.
\modelname therefore uses an emission gate.

\textbf{How.} The routed summary is emitted only if it satisfies at least one
actionability condition:
\begin{equation}
    I_t =
    \mathbb{I}
    \left[
    \exists a\in\actions_t:
    a\in\mathrm{Priority}(Z_t^\star)
    \right]
    \vee I_{\mathrm{visible}}
    \vee I_{\mathrm{held}}
    \vee I_{\mathrm{failure}}
    \vee I_{\mathrm{workflow}}.
\end{equation}
The final memory context is
\begin{equation}
    \mathcal{C}_t =
    \begin{cases}
    \mathrm{Render}(Z_t^\star), & I_t=1,\\
    \varnothing, & I_t=0.
    \end{cases}
\end{equation}
Failure-only memory is emitted only when it also yields a supported alternative
source action.

\textbf{Collaboration role.} The gate prevents collaborators' memories from
becoming passive prompt clutter. Shared experience is allowed to influence the
agent only when it can support a concrete current decision or a clearly relevant
workflow cue.

\subsection{Step 9: Render a Memory Decision Summary}

\textbf{Why.} LLM agents are sensitive to prompt length and instruction conflict.
If memory is injected as long trajectories or many retrieved snippets, it can
distract from the environment observation. \modelname therefore renders memory as
a fixed, concise decision summary.

\textbf{How.} The emitted memory block has seven fields:
\begin{quote}
\small
\begin{tabular}{@{}ll@{}}
\toprule
Field & Purpose \\
\midrule
Current phase & Identifies the macro decision phase. \\
Current state & Records target, tool, destination, held, visible, and searched state. \\
Local memory & Provides current-state graph grounding. \\
Global memory & Provides transferable workflow or source-role guidance. \\
Failure memory & Warns about similar failure patterns. \\
Next priority & Gives one concrete priority or a safe fallback. \\
Confidence / caveat & Explains whether memory is grounded or only advisory. \\
\bottomrule
\end{tabular}
\end{quote}

\textbf{Collaboration role.} The rendered summary is the final interface between
collaborative memory and individual decision making. It turns distributed
experience into a compact, phase-aware, action-oriented prompt signal while
preserving the agent's original autonomy.

\subsection{Optional TextGrad Optimization}

\textbf{Why.} Some cases require careful wording: local memory may be concrete,
global memory may be advisory, and failure memory may be useful only as a caveat.
To improve the prompt form without changing the environment workflow, \modelname
optionally optimizes only the memory summary text.

\textbf{How.} Given routed evidence $Z_t^\star$, \modelname can solve
\begin{equation}
    \mathcal{C}^{\star}_t =
    \arg\min_{\mathcal{C}\in\Omega}
    J_{\psi}
    \left(
    \mathcal{C};
    q_t,Z_t^\star,\actions_t
    \right),
\end{equation}
where $\Omega$ constrains the output to the seven-field format, forbids concrete
global location transfer, forbids new source search after the target is held, and
requires local recommendations to map to current admissible actions. If the
optimization is unavailable or unnecessary, \modelname uses the deterministic
routed summary.

\textbf{Collaboration role.} This step improves how collaborative memory is
communicated, without giving the memory module authority to replace the solver or
override environment feedback.

\end{document}
